\par\noindent\SourceLeaf{59}%
sive brevius
\[
= \frac{1}{2^{2h}}\,{}^{\nu}\mathfrak{B}^{(2h)} .  {}^{2h}\mathfrak{B}^{(h)}\,\frac{\nu+2}{h+1}\,n^{2h+1}.
\]
Sumatur porro, posita $n$ variabili, $\int An\,\mathrm{d}n$, cuius terminus initialis $=\frac{1}{2}nn$, terminus generalis $h^{\text{tus}}$
\[
= \frac{1}{2^{2h}}\,{}^{\nu}\mathfrak{B}^{(2h)} .  {}^{2h}\mathfrak{B}^{(h)}\,\frac{1}{2h+2}\,n^{2h+2}.
\]
Quae nova series, si per
\[
\frac{2(\nu+2)}{n}
\]
multiplicetur, dabit initialem $=(\nu+2)n$, generalem $h^{\text{tum}}$
\[
= \frac{1}{2^{2h}}\,{}^{\nu}\mathfrak{B}^{(2h)} .  {}^{2h}\mathfrak{B}^{(h)}\,\frac{\nu+2}{h+1}\,n^{2h+1}.
\]
Consentientibus igitur ab utraque parte terminis, fit
\[
2n .  A + A^{(1)} = \frac{2 . (\nu+2)}{n}\int An\,\mathrm{d}n,
\]
sumto sine addita constante integrali. Fluit inde, esse
\[
A^{(1)} = \frac{2 . (\nu+2)}{n}\int An\,\mathrm{d}n - 2n .  A.
\]
\noindent\hspace*{\parindent}Hac igitur formula $A^{(1)}$ ad inventam $A$ reducitur, atque haec sola restat ex genesi seriei ipsius, ut antea factum est, derivanda.

8. Confectis quae in problemate generalissime proposito sperari poterant, superest, ut ad casus speciales animum advertamus; in illis enim contrahere formulas universales nonnunquam licet. Cuius rei exemplum praebet casus, quo est $\nu$ numerus integer positivus. Tunc enim quaelibet serierum, antea ad coefficientes ipsius $(1+n\cos\varphi)^{\nu}$ exprimendos inventarum, cum coefficientes binomiales ad exponentem $\nu$ pertinentes, suo ordine quilibet, contineant, abrumpatur necesse est, et finita evadit pro quavis expressio. Simili ratione eo casu, quo $\nu$ est numerus integer negativus, ad expressiones finitas singulorum coefficientium pervenire licet. Cum ad hunc casum formulae usus plurimum revocetur, e re erit in illo enucleando paulum morari. Artificium, quo in hac disquisitione opus est, duabus absolvitur partibus; prima, ut pro coefficientibus, existente $\nu=-1$, eruantur expressiones finitae; altera, ut via monstretur, qua, concessis ipsius $(1+n\cos\varphi)^{\nu}$ coefficientibus, potestatis uno gradu inferioris
